\section{Performance Simulation}

\begin{multicols}{2}
Evaluating the definitive performance potential of the 2026 Formula 1 architecture requires mathematically rigorous, time-resolved simulation. This section presents original computational results generated via a custom MATLAB solver utilizing a Quasi-Steady-State Point Mass Model on the Autodromo Nazionale di Monza (5.793 km) \cite{monza_circuit_data}. The equations herein represent the fundamental physics governing the numeric integration.

\subsection{Simulation Methodology}

The foundational tier of the simulation establishes a discretized point mass model \cite{laptimesim_theory}. The vehicle is statically treated as an 800 kg point mass ($m$), while the Monza circuit is geometrically discretized into ultra-high-resolution segments representing straights, braking zones, and corners. 

At each integration step, the MATLAB solver calculates the absolute physical limit of the vehicle. In a purely longitudinal acceleration zone, the traction-limited acceleration ($a$) is bounded by the available hybridized powertrain output:

\begin{equation} \label{eq:sim_traction}
a = \frac{P_{ICE} + P_{elec}}{m \cdot v}
\end{equation}

where $P_{ICE}$ and $P_{elec}$ are the instantaneous power outputs of the combustion and electrical units, and $v$ is the longitudinal velocity. In a steady-state cornering segment, the maximum navigable velocity ($v_{max}$) is governed by the localized corner radius ($R$) and total available grip:

\begin{equation} \label{eq:sim_cornering}
v_{max} = \sqrt{\mu g R + \left( \frac{F_{aero}}{m} \right) \mu R}
\end{equation}

where $\mu$ is the effective tire friction coefficient, $g$ is gravitational acceleration, and $F_{aero}$ is aerodynamic downforce. During severe deceleration, the braking limit is modeled empirically, resolving to a peak deceleration of $a_{brake} = -4.5g$. 

Crucially, aerodynamic forces are continuously computed per timestep dependent on the instantaneous velocity squared:

\begin{equation} \label{eq:sim_aero}
\begin{split}
F_{down} &= 0.5 \rho C_l A v^2 \\
F_{drag} &= 0.5 \rho C_d A v^2
\end{split}
\end{equation}

where $\rho$ is the air density, $A$ is frontal area, $C_l$ is coefficient of lift, and $C_d$ is coefficient of drag. The state is numerically integrated at $dt = 0.01\mathrm{s}$.

\subsection{GGG Friction Ellipse}

Combining the absolute longitudinal and lateral limits synthesizes the dynamic operating envelope, visualized as the G-G-G Friction Ellipse \cite{pacejka_tire_lit}. The acceleration vector is constrained by:

\begin{equation} \label{eq:sim_ellipse}
\left( \frac{a_x}{a_{x,max}} \right)^2 + \left( \frac{a_y}{a_{y,max}} \right)^2 \le 1
\end{equation}

where $a_x$ is longitudinal acceleration relative to its limit $a_{x,max}$, and $a_y$ is the lateral equivalent. Aerodynamic downforce dramatically expands this performance envelope at speed. At 252 km/h (simulated exit velocity of the chicane), the car generates a $\sim$3.5g envelope. However, at 360 km/h (full straight), the immense downforce artificially expands this envelope to $\sim$5.5g. Because Monza is fundamentally a low-downforce, high-speed circuit, it predominately exploits longitudinal acceleration events, pushing the ellipse to its traction and braking boundaries.
\end{multicols}

\begin{figure}[H]
\centering
\includegraphics[width=\textwidth]{"figure13_friction_ellipse.png"}
\caption{GGG friction ellipse at 252 km/h and 360 km/h (MATLAB simulation).}
\label{fig:sim_ggg_ellipse}
\end{figure}

\begin{multicols}{2}
\subsection{Energy Deployment Simulation}

To optimize lap time while respecting the physical state of charge (SOC) degradation limit, the MATLAB solver implements a tiered State of Charge deployment logic with a hard energy budget gate. Deployment is only permitted while:

\begin{equation} \label{eq:sim_budget_gate}
E_{deployed} < E_{harvested} + 3.0\,\mathrm{MJ}
\end{equation}

To maximize the immense 350 kW Motor Generator Unit-Kinetic (MGU-K) potential, the tiers act dynamically:
1. \textbf{SOC $>$ 60\%}: Full 350 kW permitted up to 4.0s per straight.
2. \textbf{SOC $>$ 40\%}: 210 kW authorized up to 3.0s.
3. \textbf{SOC $>$ 30\%}: 105 kW restricted to 2.0s.
4. \textbf{Corner Exit}: 140 kW deployment up to 1.5s to supplement ICE torque.

The macroscopic energy balance equation strictly enforces neutrality:

\begin{equation} \label{eq:sim_energy_balance}
E_{deployed} = E_{harvested} + E_{ICE\_supplement}
\end{equation}

The actual results over the 5.793 km lap yielded an active $E_{deployed} = 4.005$ MJ. This was reconciled perfectly against heavy braking recovery of $E_{harvested} = 1.911$ MJ, supplemented by the ICE acting as a continuous generator providing $E_{ICE\_supplement} = 2.094$ MJ.
\end{multicols}

\begin{figure}[H]
\centering
\includegraphics[width=\textwidth]{"figure14_soc_trace.png"}
\caption{Battery SOC trace over lap distance with 20\%/95\% limits marked (MATLAB simulation).}
\label{fig:sim_soc}
\end{figure}

\begin{multicols}{2}
\subsection{Lap Time Simulation Results — Monza}

Aggregating the point mass matrices yields profound simulation results. The theoretical lap time crossed at \textbf{76.79 seconds (1:16.79)}. 

Contextually, the current absolute lap record at Monza stands at 79.119s (1:19.119), set by Charles Leclerc during 2019 qualifying. The simulated 2026 vehicle projects an immense $\sim$2.3-second improvement. This blistering pace is entirely consistent with the violent capability of the 350 kW MGU-K vastly out-accelerating the current 120 kW architecture out of chicanes.

The simulation logged a maximum speed of 348.2 km/h, securing an astonishing average speed of 250.2 km/h. Dynamically, the battery SOC depleted to a minimum of 25.3\%—never hitting the 20\% physical degradation floor—thereby completely validating the tiered energy management logic. The lap concluded with a final SOC of 42.6\%, securing a healthy margin for subsequent running.

The speed trace explicitly maps this aerodynamic reliance: the long straights readily reached $340+$ km/h, while chicanes plunged the minimums to $65-80$ km/h, concluding with the Parabolica exit rocketing at $\sim$250 km/h feeding back down the main straight.
\end{multicols}

\begin{figure}[H]
\centering
\includegraphics[width=\textwidth]{"figure15_speed_trace.png"}
\caption{Speed trace over lap distance --- Monza (MATLAB simulation).}
\label{fig:sim_speed}
\end{figure}

\begin{figure}[H]
\centering
\includegraphics[width=\textwidth]{"figure16_mguk_map.png"}
\caption{MGU-K deployment and harvesting map over lap distance (MATLAB simulation).}
\label{fig:sim_mgu}
\end{figure}

\begin{multicols}{2}
\subsection{Aerodynamic Efficiency Simulation}

The aerodynamic maps utilized were derived from Reynolds-Averaged Navier-Stokes (RANS) Computational Fluid Dynamics (CFD), utilizing heavily discretized meshes in ANSYS Fluent and OpenFOAM environments.

To achieve 348.2 km/h terminal velocities, the vehicle was mapped into its ultra-low drag \textit{X-Mode} configuration at Monza. Numerical CFD outputs defined an overall $C_{l} = 2.8$ and a slippery $C_{d} = 0.65$, yielding an aerodynamic efficiency ratio ($C_l/C_d$) of 4.31. This remarkably low $C_d$ directly enables the 348 km/h top speed.

Ground effect sensitivity modeling simulated $C_l$ vs ride height, showcasing severe non-linear suction increases directly below 32 mm, which mathematically justifies the strict 30 mm rear ride height kinetic target. Furthermore, the simulated dynamic front/rear downforce split successfully maintained a strict 35/65 distribution across the entire speed range.

\subsection{Vehicle Dynamics Simulation}

Vehicle stability was directly validated by the dynamics equations. The definitive understeer gradient was derived directly from the physical static 45/55 weight distribution combined with the disparate 1600/1620 mm track widths.

Engineers meticulously targeted a slightly positive gradient, inducing a highly stable, mild understeer characteristic at the absolute threshold of grip to ensure driver confidence through ultra-high-speed corners. Furthermore, implementing the ultra-low 270 mm center of gravity (CoG) directly into the model confirmed that longitudinal and lateral load transfer phenomena are highly minimized, keeping both the front and rear axles perpetually within their linear tire operating ranges throughout Parabolica and Curva Grande.
\end{multicols}

\begin{table}[H]
\centering
\caption{2026 Vehicle Performance Simulation Results (MATLAB Point Mass at Monza)}
\vspace{0.2cm}
\begin{tabularx}{\textwidth}{@{}Xl@{}}
\toprule
\textbf{Parameter} & \textbf{Value} \\
\midrule
Circuit & Autodromo Nazionale di Monza \\
Circuit Length & 5.793 km \\
Simulated Lap Time & 76.79s (1:16.79) \\
Current Lap Record (2019) & 79.119s (1:19.119) \\
Projected Improvement & $\sim$2.3s \\
Max Speed & 348.2 km/h \\
Average Speed & 250.2 km/h \\
Energy Deployed & 4.005 MJ \\
Energy Harvested & 1.911 MJ \\
ICE Supplement & 2.094 MJ \\
Min SOC & 25.3\% \\
Final SOC & 42.6\% \\
Simulation Timestep & 0.01s \\
Vehicle Model & Point Mass, 800kg \\
Aero Model & Per-timestep RANS-derived forces \\
\bottomrule
\end{tabularx}
\label{tab:monza_results}
\end{table}

